\documentclass[letterpaper,10pt]{article}

\usepackage[T1]{fontenc}
\usepackage{charter}

\usepackage{latexsym}
\usepackage[empty]{fullpage}
\usepackage{titlesec}
\usepackage{marvosym}
\usepackage[usenames,dvipsnames]{color}
\usepackage{verbatim}
\usepackage{enumitem}
\usepackage[hidelinks]{hyperref}
\usepackage{fancyhdr}
\usepackage[english]{babel}
\usepackage{tabularx}
\input{glyphtounicode}

\pagestyle{fancy}
\fancyhf{}
\fancyfoot{}
\renewcommand{\headrulewidth}{0pt}
\renewcommand{\footrulewidth}{0pt}

\addtolength{\oddsidemargin}{-0.5in}
\addtolength{\evensidemargin}{-0.5in}
\addtolength{\textwidth}{1in}
\addtolength{\topmargin}{-.7in}
\addtolength{\textheight}{1.2in}
\setlength{\footskip}{4.08003pt}

\urlstyle{same}
\raggedbottom
\raggedright
\setlength{\tabcolsep}{0in}

% Plain bold headers instead of small-caps -- small caps rendering is what
% some ATS parsers misread or skip entirely.
\titleformat{\section}{
  \vspace{-4pt}\bfseries\raggedright\large
}{}{0em}{}[\color{black}\titlerule \vspace{-5pt}]

\pdfgentounicode=1

\newcommand{\resumeItem}[1]{
  \item\small{
    {#1 \vspace{-2pt}}
  }
}

% Clean single-column, left-aligned heading layout -- no tabular/table
% environment and no right-aligned date column, since some ATS parsers
% misread a two-column layout as a table and reorder or merge the text.
% #1 Organization/Institution, #2 Location, #3 Title/Degree, #4 Dates
\newcommand{\resumeSubheading}[4]{
  \vspace{-2pt}\item
    \textbf{#1} -- #2\\
    \textit{\small#3} -- \small#4\\
    \vspace{-7pt}
}

\newcommand{\resumeSubSubheading}[2]{
    \item
    \textit{\small#1} -- \small#2\\
    \vspace{-7pt}
}

% #1 Project/activity name and stack (pre-formatted), #2 Dates
\newcommand{\resumeProjectHeading}[2]{
    \item
    \small#1 -- \small#2\\
    \vspace{-7pt}
}

\newcommand{\resumeSubItem}[1]{\resumeItem{#1}\vspace{-4pt}}

\renewcommand\labelitemii{$\vcenter{\hbox{\tiny$\bullet$}}$}

\newcommand{\resumeSubHeadingListStart}{\begin{itemize}[leftmargin=0.15in, label={}]}
\newcommand{\resumeSubHeadingListEnd}{\end{itemize}}
\newcommand{\resumeItemListStart}{\begin{itemize}[itemsep=0pt, parsep=0pt]}
\newcommand{\resumeItemListEnd}{\end{itemize}\vspace{-5pt}}

%-------------------------------------------
\begin{document}

%----------HEADING----------
% Name kept plain bold (no small caps) -- avoids styled-font parsing issues.
% Pipe separators replaced with dashes -- some ATS parsers misread "$|$".
\begin{center}
    \textbf{\fontsize{16}{18}\selectfont Daniel Larson} \\ \vspace{1pt}
    \small Santa Clarita, CA -- 661-476-0817 -- \href{mailto:lanieldarson@gmail.com}{\underline{lanieldarson@gmail.com}} --
    \href{https://www.linkedin.com/in/daniellarson1}{\underline{linkedin.com/in/daniellarson1}} --
    \href{https://github.com/laniel123}{\underline{github.com/laniel123}}
\end{center}

%-----------SUMMARY-----------
\section{Summary}
\small Builds end to end ML pipelines, NLP systems, and forecasting models that turn raw data into measurable outcomes, with hands on experience across sports analytics and startup settings using Python, R, and SQL.\vspace{-8pt}

%-----------EDUCATION-----------
\section{Education}
  \resumeSubHeadingListStart
    \resumeSubheading
      {University of California, Santa Barbara}{Santa Barbara, CA}
      {B.S. in Statistics and Data Science, GPA: 3.50, Dean's Honors}{Sep. 2024 -- Dec. 2026}
      \resumeItemListStart
        \resumeItem{Relevant Courses: Principles of Data Science, Regression Analysis,
        Applied Stochastic Processes, Statistical Learning, Design and Analysis of Experiments,
        Time Series, Big Data Analysis, STAT Data Science.}
      \resumeItemListEnd
      \resumeSubheading
      {College of the Canyons}{Santa Clarita, CA}
      {General Education, GPA: 3.80}{Aug. 2021 -- Jun. 2024}
      \resumeSubHeadingListEnd

%-----------EXPERIENCE-----------
\section{Experience}
  \resumeSubHeadingListStart

      \resumeSubheading
      {Shanghai Buba Sports \& Culture Co., Ltd. (Tenart)}{Shanghai, China}
      {Tennis Platform Operations Intern}{Jun. 2026 -- Aug. 2026}
      \resumeItemListStart
        \resumeItem{Built a SQLite database pipeline integrating the Feishu API with platform trend logging, backed by a custom validation framework for reproducible analysis.}
        \resumeItem{Implemented a Dianping web scraper capturing 500+ reviews across 162 Shanghai tennis venues, applying NLP clustering to translated reviews to quantify competitive gaps such as parking friction (cited 324 times across 40 venues).}
        \resumeItem{Fit an ARIMA time series model (ADF stationarity testing, ACF/PACF order selection, AIC based comparison) on 35+ weeks of Baidu Index data, linking Chinese tennis interest to major tournaments and projecting a 15\% increase around the 2026 Shanghai Masters.}
        \resumeItem{Delivered weekly bilingual (Chinese/English) KPI reports translating trend analysis and competitive benchmarks for the founder, while balancing on court coaching with data science responsibilities.}
      \resumeItemListEnd

  \resumeSubHeadingListEnd

%-----------PROJECTS-----------
\section{Projects}
    \resumeSubHeadingListStart

      \resumeProjectHeading
      {\textbf{SpiceRack} -- \emph{Python, Flask, Scikit-learn, NLP, TF-IDF, K-Means, SQLite, REST APIs}}{Oct. 2025 -- May. 2026}
      \resumeItemListStart
       \resumeItem{Led development as project manager for a machine learning based recipe recommendation web app spanning ML, backend, and frontend tracks.}
       \resumeItem{Designed a dual matrix NLP pipeline (TF IDF, TruncatedSVD, MiniBatchKMeans) trained on 2.2M+ recipes across 475 flavor clusters, resolving a catch all cluster failure mode validated through manual review.}
       \resumeItem{Cut recommendation latency and enabled seamless AJAX interactions by engineering a top K cosine similarity search with thread safe REST API image caching and a SQLite backed saved recipes system.}
       \resumeItem{Placed 3rd out of 100+ teams at the UCSB Data Science Club Project Showcase.}
      \resumeItemListEnd

      \resumeProjectHeading
        {\textbf{MLB Pitcher Injury Risk Prediction} -- \emph{R, XGBoost, K-Means, tidymodels}}{Jan. 2026 -- Mar. 2026}
        \resumeItemListStart
          \resumeItem{Predicted MLB pitchers' injury risk for the previous season by building and validating 4 ML models (logistic regression, elastic net, random forest, XGBoost) on 30 years of Lahman Database pitching data, identifying a player's total innings pitched as the leading risk factor.}
          \resumeItem{Revealed workload as a key injury driver by segmenting pitchers into 3 archetypes via K-means clustering; the highest workload group carried nearly double the injury rate (26.6\%) of the lowest (14.1\%).}
          \resumeItem{Tuned the final model to correctly flag 80\% of at risk pitchers, with the top 10\% highest risk group injured at a rate of 53.6\%, which is nearly 3x the 19\% baseline average.}
        \resumeItemListEnd

    \resumeSubHeadingListEnd

%-----------ACTIVITIES-----------
\section{Activities \& Leadership}
    \resumeSubHeadingListStart
      \resumeProjectHeading
          {\textbf{Data Science Club (UCSB)} -- \emph{Member}}{Oct. 2024 -- Present}
          \resumeItemListStart
          \resumeItem{Won Best Sports Analyst Project at Datathon 2025 and reached the Project Pitchfire final twice among 100+ competing teams.}
          \resumeItemListEnd
    \resumeSubHeadingListEnd

%-----------SKILLS-----------
\section{Skills}
 \begin{itemize}[leftmargin=0.15in, label={}]
    \small{\item{
     \textbf{Programming Languages}{: Python, R, SQL, C++} \\
     \textbf{Concepts}{: Machine Learning, Statistical Modeling, Classification, Clustering, Regression, Natural Language Processing, Tokenization, Time Series Forecasting, Feature Engineering, A/B Testing, REST APIs} \\
     \textbf{Libraries \& Frameworks}{: Scikit-Learn, XGBoost, tidymodels, statsmodels, Flask, Pandas, NumPy, Selenium, BeautifulSoup} \\
     \textbf{Tools}{: Git, GitHub, Jupyter Notebook, VS Code, Docker, SQLite, Feishu} \\
     \textbf{Languages (Spoken)}{: English (Native), Spanish (Elementary), Mandarin Chinese (Beginner)}    }}
 \end{itemize}

\end{document}
